\documentclass[11pt]{article}
\usepackage[a4paper,margin=25mm]{geometry}
\usepackage{amsmath,amssymb,amsthm}
\newtheorem{theorem}{Theorem}
\newtheorem{lemma}{Lemma}
\title{Structural certificates for the $\mathrm{Sp}(6,2)$ atlas table}
\author{L054 finite application}\date{}
\begin{document}\maketitle

\begin{theorem}
Relative to the acquired numerical catalogue and named structural registry,
all $1369$ rows of the current companion table headed $\mathrm{Sp}(6,2)$
have the displayed structural descriptions. The ambient group is
$\operatorname{PSp}_6(2)$, of order $1451520$, in the acquired action on
$63$ points. An extension description asserts the indicated normal
subgroup, quotient and splitting behaviour. It does not determine an
abstract isomorphism class uniquely from the printed expression.
\end{theorem}

\paragraph{Source and acquired results.}
The source is the Connor--Leemans companion table dated 7 June 2013,
associated with \cite{atlas}. Its numerical catalogue, source-row
bijection, actual embedded subgroup generators and named top group are
acquired from the separately accepted numerical computation. No new
ambient subgroup census is used here. The $515$ acquired structural models
come from earlier accepted finite applications. Each of the $52$ files
read to reconstruct this registry agrees with its entry in the accepted
$\mathrm{Sp}(4,7)$ review's dependency record. Acquiring these inputs
retains their existing hypotheses and qualifications.

The paper's Section 3.2 adopts left association for an unparenthesized
sequence of extensions. We have read that convention on native printed
page 263. A position audit joins the powers, subscripts, operators,
parentheses and wrapped structural cells of all $1369$ rows to their parsed
expressions and orders. Direct inspection of the companion covers pages
$1,2,49,86,87$. It is not an inspection of every page. The position audit
uses the same pypdf library as the flat extraction, so these are not two
independent extraction engines. The exact retained files, rather than an
assertion of identity with every historical edition, specify the source.

\paragraph{What is supplied by this application.}
There are $199$ new finite model certificates: $143$ split extensions,
$35$ direct products and $21$ nonsplit extensions. Every proper source
occurrence has a transport from a proved model to its acquired subgroup
of the $63$-point action. The source row, its literal expression, its order
and its actual generators are joined before the transport is checked.
Candidate models used to help later searches are still proved individually;
they are never promoted to acquired facts merely because a search returned
them. All references to new factor models precede the model using them.

\begin{lemma}
Let $M$ and $H$ be finite groups given by complete closures of specified
permutation generators. A bijection $f:M\to H$ satisfying $f(1)=1$ and
$f(xg)=f(x)f(g)$ for every $x\in M$ and every specified generator $g$ is
an isomorphism.
\end{lemma}
\begin{proof}
Every element $y\in M$ is a positive word in the generators: in a finite
group, an inverse of a generator is a nonnegative power of that generator.
Induction on such a word gives $f(xy)=f(x)f(y)$ for every $x,y\in M$.
Thus $f$ is a homomorphism, and the assumed bijection makes it an
isomorphism. The reader constructs the proposed images along the complete
closure, checks their consistency whenever a word reaches an old element,
and then checks every right-generator edge. Bijectivity is checked against
the whole target subgroup, not merely against its generating set.
\end{proof}

\paragraph{Extension certificates.}
For each new model $H$, the reader constructs the subgroup $N$ generated by
the displayed permutations, verifies $N\triangleleft H$, and transports an
earlier model onto $N$. For a split extension it similarly verifies a
subgroup $K$, its structure, $N\cap K=1$ and $|N||K|=|H|$. Normality makes
$NK$ a subgroup, and the order equation gives $H=NK$. This proves the
semidirect product assertion. For a colon the retained convention also
requires nontrivial actual conjugation of $K$ on $N$; a pair of
noncommuting generators supplies the witness. For a direct product all
generator pairs commute, hence all elements of the two factors commute.
These are the previously accepted finite extension predicates.

\paragraph{The nonsplitting certificates.}
Thirteen models have prime quotient. If that quotient has order $p$, a
complement would be generated by an element of order $p$ outside $N$.
Conversely such an element generates a complement by the order formula.
The reader checks every outside element and excludes this possibility.
One model has an abelian quotient. Its certificate gives complete,
disjoint normal forms for the quotient and checks every choice of lifts
of its stated cyclic generators. Every choice fails an order relation
or a commutation relation, so none defines a section.

Five models have quotient $S_3$. The certificate identifies quotient
generators satisfying $a^2=b^3=(ab)^2=1$ and gives all their lifts. A
complement would provide lifts satisfying these same relations. First
discarding lifts which fail $a^2=1$ or $b^3=1$ preserves every possible
section; the remaining pairs are checked against $(ab)^2=1$. The source
count records the full original Cartesian product, including discarded
pairs. This filtering is inherited from the earlier accepted reader.

The remaining two dots use the previously accepted Cayley-graph criterion.
The quotient is identified with an acquired finite model. For each choice
of lifts of its ordered generators, the reader propagates a proposed
section along the complete quotient Cayley graph. Consistency on every
edge would give a homomorphism by the lemma, and the quotient images would
make it a section. Conversely a section supplies such a consistent choice.
One inconsistent edge for every lift choice therefore proves that no
complement exists. The present instances require $128$ lift choices and
$2176$ edge tests. A failed GAP complement search is never a certificate
of nonsplitting.

\paragraph{The last source expression.}
Row 28 has order $1536$ and expression
\[
  2^2\times4\mathbin{\cdot}2:2:2:2:S_3.
\]
Left association gives successive orders
$16,32,64,128,256,1536$. In particular, the direct product at the start
is formed before the nonsplit extension. The last four newly generated
models complete the required chain from order $64$ onwards and its final
$S_3$ complement. The literal reader proves these models by the same
predicates as every other row; no special interpretation is introduced
to obtain this result.

The proposal search used two elementary improvements. For an abelian
requested expression it rejects a subgroup with different cyclic
prime-power invariants, after checking abelianness. These invariants are
obtained by multiset union for direct products. It also requires the
requested action before selecting the first prime-order complement
generator. An earlier selection could choose a centralizing involution
and overlook a later involution with nontrivial action. Both changes
concern proposal efficiency and completeness, and leave the proof
predicates unchanged. They do not constitute new structure theorems.

\paragraph{Verification and attribution.}
The complete literal check passed in $33.710$ seconds. It checked all
$1369$ source joins, the $199$ new model proofs, $1747$ transports and
$1682351$ generator edges. The nonsplitting checks include $1376$ outside
elements and $9472$ composite-quotient lift pairs, in addition to the
Cayley-graph counts above. Input hashes remained unchanged and stderr was
empty. Twenty damaged source or certificate inputs are rejected in both
ordinary and optimized Python. The complete optimized replay reproduces
the same proof counts. Source admission is recorded separately by the
distinct review, with the extraction qualifications stated above.

The original table, its notation and association convention are due to
the source authors. The numerical catalogue and named models are acquired
results. The finite transport, extension and lift criteria are reused
from the project's accepted common method. The contribution here is their
explicit application to every occurrence of this table, with literal
certificates and an exact source join. No new conjecture resolution,
Lean formalization, historical-byte identity, whole-article completion or
compression measurement is claimed. The other large source tables remain
in scope.

\begin{thebibliography}{1}
\bibitem{atlas} T. Connor and D. Leemans,
An atlas of subgroup lattices of finite almost simple groups,
\emph{Ars Mathematica Contemporanea} \textbf{8} (2015), 259--266.
The retained companion table is headed $\mathrm{Sp}(6,2)$ and dated
7 June 2013.
\end{thebibliography}
\end{document}
